% Original teaching diagram from the exact example in Chapter 19.
\begin{tikzpicture}[x=1mm,y=1mm,font=\figurefont\fontsize{8}{10}\selectfont,
 v/.style={circle,draw=FigureStroke,fill=white,minimum size=6mm,inner sep=0pt},
 w/.style={fill=white,inner sep=1pt},
 sub/.style={font=\FigureNoteFont\fontsize{7.5}{9}\selectfont,text=FigureInk}]
  \foreach \shift/\i in {0/a,51/b} {
    \begin{scope}[xshift=\shift mm]
      \coordinate (p) at (0,2);
      \coordinate (q) at (17,25);
      \coordinate (r) at (34,2);
      \draw[FigureStroke] (p)--(r) node[midway,w,below=1pt] {$1$};
      \ifnum\shift=0
        \draw[FigureStroke] (p)--(q) node[midway,w,left=1pt] {$1$};
        \draw[FigureStroke] (q)--(r) node[midway,w,right=1pt] {$1$};
      \else
        \draw[FigureBlue,line width=1.35pt] (p)--(q) node[midway,w,left=1pt] {$3/2$};
        \draw[FigureRed,densely dashed] (q)--(r) node[midway,w,right=1pt] {$1/2$};
      \fi
      \node[v] at (p) {$1$};
      \node[v] at (q) {$2$};
      \node[v] at (r) {$3$};
      \node[below=3.5mm] at (p) {$x_1=3$};
      \node[above=4mm] at (q) {$x_2=-2$};
      \node[below=3.5mm] at (r) {$x_3=-1$};
      \node[sub] at (17,-11) {(\i)};
    \end{scope}
  }
\end{tikzpicture}
